%%%%%%%%%%%%%%%%%%%%%%%%%%%%%%%%%%%%%%%%%%%%%%%%%%%%%%%
% SECTION 9: DISCUSSION
%%%%%%%%%%%%%%%%%%%%%%%%%%%%%%%%%%%%%%%%%%%%%%%%%%%%%%%
% 
% INPUTS NEEDED: 
%   - Synthesis of all results
%   - Comparison to general LLMs
%   - Error analysis summary
%   - Future directions
%
% FIGURES: None
% TABLES: None
% DEPENDENCIES: All Results sections (references metrics)
%
%%%%%%%%%%%%%%%%%%%%%%%%%%%%%%%%%%%%%%%%%%%%%%%%%%%%%%%

\section{Discussion}\label{sec3}

%%% PARAGRAPH 1: System readiness + modular architecture %%%
OncoCITE is a viable system ready for institutional testing in clinical settings. The modular architecture facilitates integration with existing hospital information systems, while comprehensive logging and transparency features address regulatory requirements for clinical-decision-support tools. While large general-purpose LLMs such as GPT-4 or Claude can answer oncology questions, their responses are limited by a static knowledge cutoff, occasional hallucinated citations, and an absence of structured, auditable output. By contrast, OncoCITE couples an LLM to a retrieval layer, canonical dictionaries, and explicit ReACT traces, thereby providing source-linked, CIViC-compatible evidence items that satisfy traceability requirements.

%%% PARAGRAPH 2: Terminology challenges + canonicalization solution %%%
Clinicians rarely use standardized terminology; they speak in abbreviations, synonyms, and colloquialisms shaped by training and institutional culture. A surgeon may document ``mammary carcinoma'' whereas an oncologist notes ``breast cancer''; both refer to the same disease but can confound systems trained on rigid terminology. Similarly, a single gene may be referred to by multiple aliases (for example, ERBB2/HER2), and therapy nomenclature varies, where ``trastuzumab'' and ``Herceptin'' denote the same monoclonal antibody, whereas ``anti-HER2 therapy'' is a broader class label. Early development versions of OncoCITE failed silently on such variants. Incorporating canonicalization dictionaries and iterative query reformulation eliminated these silent failures, rescuing most previously unsuccessful requests and markedly raising overall query success without appreciably lengthening the 6--7~second response time.

%%% PARAGRAPH 3: Clinical Interpretation Engine %%%
The Clinical Interpretation Engine bridges the gap between accurate retrieval and actionable insight. Rather than returning raw tables, it generates structured, hierarchical reports: concise summaries for immediate context, key findings highlighting therapeutic or prognostic significance, relevance to current treatment paradigms, and suggested next steps. Simulated GPT-4.1 assessments showed that queries containing typographical errors or colloquial terminology, previously total failures, were transformed into clinically contextualized answers. Representative test cases confirmed that challenging questions, which confound conventional interfaces, are now resolved by OncoCITE.

%%% PARAGRAPH 4: Error tracking + safety %%%
Comprehensive tracking across all phases revealed distinct failure patterns. Most extraction errors involved complex variant nomenclature, whereas natural-language-to-SQL performance dipped with non-standard phrasing until the knowledge layer intervened. Clarification prompts, informative failure messages, and a conservative uncertainty-management policy prevented potentially harmful recommendations; no unsafe advice was observed. Audit trails record every canonical mapping, constraint relaxation, and SQL reformulation, providing provenance essential for clinical validation.

%%% PARAGRAPH 5: ReACT performance validation %%%
The performance achieved by the ReACT-enhanced extraction pipeline demonstrated that transparent automation can match human quality at ten-fold the speed. The GPRC5D case study further illustrated real-world utility, automatically extracting five clinical trials for emerging bispecific and CAR-T therapies with 100\% concordance to published response rates and generating CIViC-compatible evidence items absent from the database. This work establishes OncoCITE as a data-extraction and access pipeline. A separate, ongoing study will develop and evaluate patient-level report generation and treatment prioritization in prospective settings.

%%% PARAGRAPH 6: Broader impact + applications %%%
The successful deployment of OncoCITE could fundamentally transform how clinicians interact with genomic knowledge bases, potentially increasing the clinical utility of genetic testing and improving patient outcomes through more informed therapeutic decisions. As an extraction and access layer, OncoCITE enables standards-aligned, auditable evidence to flow into downstream applications; patient-level reporting and treatment prioritization are beyond the scope of this work and are being developed in an ongoing study. The framework's modular architecture enables adaptation to other medical domains where structured knowledge extraction and natural language accessibility represent critical barriers. Beyond oncology, similar approaches could enhance clinical decision-making in rare genetic diseases, pharmacogenomics, and emerging fields like liquid biopsy interpretation where the gap between research discoveries and clinical application remains substantial.

%%% PARAGRAPH 7: Transparency as educational tool %%%
The transparent reasoning chains generated by our ReACT implementation may serve as educational tools, helping clinicians understand the evidence evaluation process while maintaining trust through auditability. This transparency becomes particularly valuable as regulatory frameworks evolve to address AI in clinical decision support, where explainability and traceability are essential requirements.

%%%%%%%%%%%%%%%%%%%%%%%%%%%%%%%%%%%%%%%%%%%%%%%%%%%%%%%
% DISCUSSION STRUCTURE SUMMARY:
%
% Para 1: System readiness
%   - Institutional testing ready
%   - Modular architecture
%   - Comparison to GPT-4/Claude
%   - OncoCITE advantages: retrieval + dictionaries + ReACT traces
%
% Para 2: Terminology challenge
%   - Clinician language variability
%   - Examples: mammary carcinoma/breast cancer, ERBB2/HER2
%   - trastuzumab/Herceptin example
%   - Canonicalization solution
%   - 6-7 second response maintained
%
% Para 3: Clinical Interpretation Engine
%   - Hierarchical reports structure:
%     • Concise summaries
%     • Key findings
%     • Treatment relevance
%     • Suggested next steps
%   - Typo/colloquial handling success
%
% Para 4: Safety and error tracking
%   - Distinct failure patterns per phase
%   - Clarification prompts
%   - Uncertainty management
%   - No unsafe advice observed
%   - Complete audit trails
%
% Para 5: Validation results
%   - 10x speed vs human
%   - GPRC5D: 5 trials, 100% concordance
%   - Scope statement: extraction and access only
%   - Future work: patient-level reports (separate study)
%
% Para 6: Broader impact
%   - Transform clinician-database interaction
%   - Improve genetic testing utility
%   - Other domains: rare diseases, pharmacogenomics, liquid biopsy
%
% Para 7: Educational value
%   - Reasoning chains as learning tools
%   - Regulatory compliance (explainability)
%
%%%%%%%%%%%%%%%%%%%%%%%%%%%%%%%%%%%%%%%%%%%%%%%%%%%%%%%

%%%%%%%%%%%%%%%%%%%%%%%%%%%%%%%%%%%%%%%%%%%%%%%%%%%%%%%
% KEY CLAIMS TO VERIFY:
%
% □ "ten-fold the speed" vs human curation
% □ "five clinical trials" from GPRC5D extraction
% □ "100% concordance to published response rates"
% □ "6-7 second response time"
% □ "no unsafe advice was observed"
%
%%%%%%%%%%%%%%%%%%%%%%%%%%%%%%%%%%%%%%%%%%%%%%%%%%%%%%%

%%%%%%%%%%%%%%%%%%%%%%%%%%%%%%%%%%%%%%%%%%%%%%%%%%%%%%%
% SCOPE LIMITATIONS EXPLICITLY STATED:
%
% "patient-level reporting and treatment prioritization are 
% beyond the scope of this work and are being developed in 
% an ongoing study"
%
% This is important - defines what OncoCITE does NOT do
%%%%%%%%%%%%%%%%%%%%%%%%%%%%%%%%%%%%%%%%%%%%%%%%%%%%%%%
